\section{Exact convergence boundary of the defining Dirichlet series}
\label{sec:dirichlet-divergence}

The convergence theorem for $\Re s>1$ has a constructive converse.
This converse concerns the defining series $\sum n^{-s}$, not the
continued function $\zeta(s)$.  A short finite argument treats the
whole complementary half-plane, including its boundary, without
asking whether the imaginary part equals zero.

Write $s=\sigma+it$, and suppose that $\sigma\le1$ and
$|t|\le B$, where $B\ge1$ is an integer.  For an arbitrary starting
index $N$, set
\[
 q=2B,\qquad K=\max(N,1),\qquad
 M=qK,\qquad L=(q+1)K.
\]
Both $M$ and $L$ are at least $N$.  For each integer $M<n\le L$,
\[
 0\le\log(n/M)\le\log(1+1/q)\le1/q,
 \qquad |t\log(n/M)|\le\tfrac12.                  \tag{V1}
\]
The logarithmic inequality follows by integrating
$(1+u)^{-1}\le1$ on $[0,1/q]$.  The finite cosine series, grouped
into alternating pairs, gives
$\cos v\ge1-v^2/2\ge7/8$ for $|v|\le1/2$; we use the weaker
bound $1/2$.

Multiply the finite block by the unit complex number $M^{it}$.
By the elementary exponential addition law,
\[
 \Re\left(M^{it}\sum_{n=M+1}^{L}n^{-s}\right)
   =\sum_{n=M+1}^{L}n^{-\sigma}\cos\bigl(t\log(n/M)\bigr).
\]
Since $\sigma\le1$ and $n\ge1$,
$n^{-\sigma}=n^{-1}e^{(1-\sigma)\log n}\ge n^{-1}$.
There are exactly $K=L-M$ terms.  Therefore
\begin{equation}
 \boxed{\displaystyle
 \left|\sum_{n=M+1}^{L}n^{-s}\right|
 \ge\frac12\sum_{n=M+1}^{L}\frac1n
 \ge\frac{K}{2L}=\frac1{2(2B+1)}.}
 \label{eq:uniform-dirichlet-noncauchy}
\end{equation}
The first inequality uses only that multiplication by $M^{it}$
preserves modulus and modulus is at least real part.  Thus the
partial sums are not Cauchy, with the explicit witnesses
\[
 \varepsilon_B=\frac1{4(2B+1)},\qquad
 (M,L)=\bigl(2B\max(N,1),(2B+1)\max(N,1)\bigr).
\tag{V2}
\]
In particular one may fix the dyadic accuracy index $p_0=2B+4$:
for every $N$, the displayed block has modulus greater than
$2^{-p_0}$.  These witnesses are primitive recursive even before any
transcendental function is evaluated.

This uniform statement includes all of the usual cases.  When
$\sigma\le0$, there is also the still simpler term test:
$|n^{-s}|=n^{-\sigma}\ge1$ for every $n\ge1$.  For
$0<\sigma<1$, no lower bound for $\sigma$ or for $1-\sigma$ is
needed in \eqref{eq:uniform-dirichlet-noncauchy}; the supplied
inequality $\sigma\le1$ suffices.  On $\sigma=1$ the same blocks
work both for $t$ apart from zero and for $t=0$.  In the latter case,
the usual dyadic harmonic block has the stronger bound
$\sum_{n=K+1}^{2K}1/n\ge1/2$.  These are separately stated
consequences, not an appeal to a decidable trichotomy for arbitrary
real names.

For the arithmetical reading, choose in addition an integer
$R\ge|s|$ to bound evaluations.  On the displayed finite block,
$n^{-\sigma}\le L^R$.  Evaluate each positive amplitude and each
cosine to rational error at most
\[
 \eta=\frac1{128(2B+1)K(1+L^R)}.
\]
The finite product-and-sum error is then smaller than
$1/[16(2B+1)]$.  Argument errors in the elementary evaluations are
reduced by their already supplied exponential and logarithmic
Lipschitz bounds.  The resulting finite rational inequality
certifies a positive real part exceeding
$1/[4(2B+1)]$.  A bounded sum of these rational error checks is a
quantifier-free induction invariant.  No decision about equality of
real numbers and no search for a favorable oscillation is needed.

Together with the explicit tail estimate in $\Re s>1$, this proves
the stated exact boundary: the Dirichlet series defines zeta on the
certified right half-plane, and is constructively non-Cauchy on the
certified closed half-plane $\Re s\le1$.  It does not license an
undecidable case distinction for an arbitrary input whose position
relative to the boundary has not been supplied.
